\section{Positive quadrature and finite data}\label{sec:finite}

The spectral formulas of Section~\ref{sec:fourier} turn coefficient lists
into genuine Fourier pairs.  We now construct ten pairs of finite coefficient
columns.  Their linear
combinations will approximate the two lists required by the coupled jet
equations.  The columns are fixed; only the coefficients of the combination
vary with the Gaussian parameter.  We control the construction in two ways:
a uniform quadrature estimate bounds the error from replacing the six density
integrals by finite sums, and finite interval comparisons establish the norm
and sign bounds for the columns.  Section~\ref{sec:interpolation} then corrects
the lists to exact interpolation, and Section~\ref{sec:signs} uses these
bounds to prove the signs.

All terminating decimals in this section denote exact rational numbers.
For a real matrix, $\norm{A}$ is its maximum absolute column sum.  A column
has its $\ell^1$ norm, denoted $\norm{v}_1$.  Matrix products below are ordinary
products, without complex conjugation.

\subsection{A positive rule for the six density integrals}

Set
\[
 M=168,\qquad \mathcal F=\mathcal N\cap[1,M],\qquad N=384.
\]
There are $84$ nodes in $\mathcal F$: each of the fourteen successive
blocks of twelve positive integers contains six permitted residues.  A finite
positive-node list has $168$ coordinates, first all $c_n$ and then all $d_n$,
with $n$ increasing in each group.  We append the extra coordinates
$e=(c_0,d_0,C)$ in that order.  Thus the complete finite input space has
$171$ coordinates.

The folded spectral measure in Lemma~\ref{lem:cardinal-measure} has a density
on each interval $((j-1)/6,j/6)$ and seven atoms at $t_j=j/6$.  We leave the
atoms unchanged.  On each density interval define
\begin{equation}\label{eq:quadrature-rule}
\begin{aligned}
 \alpha_v&=\frac{v+1/2}{N},&
 \tau_{jv}&=\frac{2j-1+\cos(\pi\alpha_v)}{12},\\
 o_v&=\frac1N\left(1-2\sum_{a=1}^{N/2-1}
             \frac{\cos(2a\pi\alpha_v)}{4a^2-1}\right),
 &0&\le v<N,\quad 1\le j\le6.
\end{aligned}
\end{equation}
The replacement of the density integral is the sum of masses
$o_vd_j(\tau_{jv})/6$.  Since $0<\pi\alpha_v<\pi$, every quadrature
point lies inside its frequency interval.  This is the first Fej\'er rule
described by Waldvogel~\citep[Section~2]{Waldvogel2006}, divided by two so that its weights on
$[-1,1]$ approximate the average rather than the integral of mass two.
We prove the exactness and positivity needed for this normalization.

Put $\theta_v=\pi(v+1/2)/N$.  For $1\le l\le2N-1$, the finite identity
\[
 \sum_{v=0}^{N-1}\cos((v+1/2)t)=\frac{\sin(Nt)}{2\sin(t/2)}
\]
at $t=l\pi/N$ gives $N^{-1}\sum_v\cos(l\theta_v)=0$.  Applying the
product-to-sum formula, for $1\le l\le N-1$ and $1\le a\le N/2-1$ we obtain
\[
 \frac1N\sum_{v=0}^{N-1}\cos(l\theta_v)\cos(2a\theta_v)
 =\begin{cases}1/2,&l=2a,\\0,&l\ne2a.\end{cases}
\]
Consequently
\[
 \sum_{v=0}^{N-1}o_v\cos(l\theta_v)
 =\begin{cases}
 1,&l=0,\\
 0,&l\text{ odd},\\
 1/(1-l^2),&l>0\text{ even},
 \end{cases}
 \qquad 0\le l\le N-1.
\]
Let $C_l$ be the Chebyshev polynomial defined by
$C_l(\cos\theta)=\cos(l\theta)$. These values equal
$\frac12\int_{-1}^1C_l(x)\,\dd x$, as follows by substituting
$x=\cos\theta$. The polynomials $C_0,\ldots,C_{N-1}$ span all polynomials of degree at most
$N-1$, so the rule is exact through degree $383$.  Moreover,
\[
 \sum_{a=1}^{N/2-1}\frac{2}{4a^2-1}=1-\frac1{N-1}
\]
by telescoping.  Since every cosine is at most one, this identity and the
case $l=0$ above prove
\begin{equation}\label{eq:positive-quadrature-weights}
 o_v\ge\frac1{N(N-1)}>0,\qquad \sum_{v=0}^{N-1}o_v=1.
\end{equation}
The substitution $t=(2j-1+x)/12$ changes the average rule on $[-1,1]$
into the weights $o_v/6$ on the $j$th frequency interval.  Each interval
therefore has total weight $1/6$, and all six together have total weight one.

A tilde will denote this replacement of the density integrals.  In particular,
$\widetilde p$ and $\widetilde K_p$ are finite sums.  The functions ultimately
used in the Fourier pair remain $p$ and $K_p$ with their exact spectral
integrals; the following lemma quantifies the difference.

\begin{lemma}[Quadrature error for scaled moments]\label{lem:quadrature-error}
Let $p$ and $K_p$ be the exact direct and transformed functions for one basis
input among the $171$ coordinates $(\mathcal F,e)$.  For every integer
$a\ge0$, real $0\le m\le168$, and real $|y|\le3/2$, the replacement
\eqref{eq:quadrature-rule}, with all atoms kept exactly, changes each of
\[
 \frac{y^a}{a!}p^{(a)}(m),\qquad
 \frac{y^a}{a!}K_p^{(a)}(m)
\]
by less than $10^{-31}$ in absolute value.
\end{lemma}

\begin{proof}
We first bound each analytic integrand on a disk twice as wide as its real
integration interval.  We then use polynomial exactness and positivity to
turn one Taylor remainder bound into a quadrature error bound.  The estimates
are uniform in the derivative order.

For the $j$th density piece in \eqref{eq:folded-density}, put
$u=(2j-1)/12$ and consider $|t-u|\le1/6$.  On this disk,
$|t|\le13/12<1.1$.  The pole $-ih$ of $z$ and $\lambda$ lies outside, and
\[
 |\lambda(t)|\le\frac1{b(h-1/6)}<5,
 \qquad
 |z(t)|\le\frac{B}{h-1/6}+h=\frac{214}{35}<6.2.
\]
Inversion sends the disk centered at $u+ih$ with radius $1/6$ to the disk
centered at $(u-ih)/(u^2+h^2-1/36)$ with radius
$(1/6)/(u^2+h^2-1/36)$.  Thus
\begin{equation}\label{eq:complex-kernel-disk}
 \Ima z(t)\ge
 \frac{B(h-1/6)}{u^2+h^2-1/36}-h
 \ge-\frac{1402}{17505}>-.081.
\end{equation}
The middle expression decreases as $u$ increases, so its minimum among the
six disks occurs at $u=11/12$.

For positive indices, the moduli of the coefficients in
\eqref{eq:periodic-coefficients} are $(2,2,2,1,2,1)$.  Hence
\[
 \sum_{l=1}^6|P_l|=10,\qquad
 \sum_{l=1}^6t_l|P_l|=\frac{16}{3}.
\]
For a $c_n$ basis column, the density on the disk has modulus at most
\[
 2\pi^2e^{\pi n/6}\sum_{l=j}^6|P_l|(t_l-u+1/6)
 \le\frac{37\pi^2}{3}e^{\pi M/6}<220e^{\pi M/6}.
\]
The last sum is bounded by $16/3+(1/12)10=37/6$; the first disk gives the
largest possible added term.  For a $d_n$ column the smaller bound
$20\pi e^{\pi M/6}$ suffices.  These arguments include $n=0$.
The constant column has only atoms, so it has no quadrature error.

The direct scaled kernel factor satisfies
\[
 |e^{i\pi tm}|\frac{|i\pi ty|^a}{a!}
 \le e^{\pi M/6}e^{1.65\pi},
\]
and the transformed factor satisfies
\[
 |\lambda(t)e^{i\pi z(t)m}|\frac{|i\pi z(t)y|^a}{a!}
 \le5e^{.081\pi M}e^{9.3\pi}.
\]
Here $x^a/a!\le e^x$ for $x\ge0$ removes all dependence on $a$.
The resulting direct and transformed integrand bounds are respectively
\[
 220e^{\pi M/3+1.65\pi},\qquad
 1100e^{\pi M/6+.081\pi M+9.3\pi}.
\]
For completeness, the elementary estimates used to bound these two numbers
need no matrix or interval computation.  Polynomial division gives
\[
 0<\int_0^1\frac{x^4(1-x)^4}{1+x^2}\,\dd x=\frac{22}{7}-\pi.
\]
The positive exponential series, whose tail after degree six is bounded by
a geometric series of ratio $1/8$, gives
\[
 e<\sum_{j=0}^6\frac1{j!}+\frac1{7!(1-1/8)}
   =\frac{95901}{35280}<\frac{2719}{1000}.
\]
With $M=168$, the two exponents above are less than $182$ and $160$,
respectively, after using $\pi<22/7$.  The exact rational comparisons
\[
 220(2719/1000)^{182}<10^{83},\qquad
 1100(2719/1000)^{160}<10^{83}
\]
therefore bound both analytic integrands by $10^{83}$.  This is the
degree-six elementary route for the present allowance; it is separate from
the scalar method developed in Appendix~\ref{sec:arithmetic}.

The actual interval is $|t-u|\le1/12$, half the disk radius.  Cauchy's
coefficient estimate bounds the Taylor tail after degree $383$ there by
\[
 10^{83}\sum_{r=384}^{\infty}2^{-r}
 =2\cdot10^{83}2^{-384}.
\]
The degree-$383$ polynomial is integrated exactly.  Both integration and
quadrature are positive functionals of total mass one when the six intervals
are combined, by \eqref{eq:positive-quadrature-weights}.  Applying each to
the tail thus bounds their difference by
\[
 4\cdot10^{83}2^{-384}<10^{-31};
\]
the final comparison is the integer inequality $4\cdot10^{114}<2^{384}$.
Taking real parts cannot enlarge the error, and no atom was changed.
\end{proof}

\subsection{Masses and transformed jets}

The quadrature is now both specified and controlled.  We next express it as
matrix multiplication, so that every entry of the finite construction is
determined by the formulas already given.  The mass matrix $W$ has $6N=2304$
quadrature rows in lexicographic $(j,v)$ order, followed by the seven atom
rows in increasing $j$ order; its $171$ columns use the coordinate order
specified above.  For $t=\tau_{jv}$ and $n\in\mathcal F\cup\{0\}$, set
\begin{align*}
 W_{(j,v),c_n}
 &=-\frac{2\pi^2o_v}{6}e^{-i\pi tn}
       \sum_{l=j}^6(t_l-t)P_le^{i\pi t_ln},\\
 W_{(j,v),d_n}
 &=\frac{2\pi i o_v}{6}e^{-i\pi tn}
       \sum_{l=j}^6P_le^{i\pi t_ln},&
 W_{(j,v),C}&=0.
\end{align*}
At atom $j$, only the constant column is nonzero:
$W_{\mathrm{atom}(j),C}=e_jP_j$, where $e_0=1$ and $e_j=2$ for $j>0$.
These are exactly the folded masses of Lemma~\ref{lem:cardinal-measure}.

For the tuple $t$ of all $2311$ mass points, define the row vectors
\begin{equation}\label{eq:quadrature-moments}
\begin{aligned}
 F_a(m,y;t)&=e^{i\pi tm}\frac{(i\pi yt)^a}{a!},\\
 G_a(m,y;t)&=\lambda(t)e^{i\pi z(t)m}
                  \frac{(i\pi yz(t))^a}{a!}.
\end{aligned}
\end{equation}
For a real coefficient column $A$, the real parts of $F_aWA$ and $G_aWA$
are the order-$a$ derivatives at $m$ of its direct and transformed finite
sums, multiplied by $y^a/a!$.  This follows by differentiating each finite
exponential sum.  It fixes, in particular, the absence of conjugation in
these products.

Apply the jet projection of \eqref{eq:jet-projection} to the transformed
sums and write
\[
 (\widetilde S_{\mathcal F\mathcal F},\widetilde S_{\mathcal F e})
 =-\mathcal J_{\mathcal F}\widetilde K,
 \qquad D=\widetilde S_{\mathcal F\mathcal F}.
\]
Thus $D$ is a $168\times168$ matrix; it is distinct from the scalar node
constants $D_n$.  Explicitly, if
$a_n=\Rea(G_0(n,1)W)$ and $b_n=\Rea(G_1(n,1)W)$, the $c_n$ and $d_n$
rows of the full projected matrix are
\[
 -a_n/Q_n,\qquad (D_na_n-b_n)/Q_n,
\]
respectively.  The bounds in \eqref{eq:node-constant-bounds} ensure that
these denominators are positive.

\subsection{Target data and finite reference columns}

The ten data coordinates encode the target jets at the first five positive
nodes.  Define
\[
 (n_1,\ldots,n_5)=(1,3,4,7,9),\qquad
 v_j=e^{-k(n_j-1)},\qquad
 q(k)=(v_1,\ldots,v_5,v_1/k,\ldots,v_5/k)^t.
\]
The first five coordinates are $-T_k'(n_j)$ and the last five are
$T_k(n_j)$.  We will convert these values and slopes to cardinal coordinates
with the jet projection $\mathcal J$.
In the cardinal formula, $c_0,d_0$ determine the value and slope of the direct
profile $p$ at the origin, while $C$ multiplies $P$ and supplies the constant
term in the later rational tail comparison.  These three coordinates leave
all direct positive-node jets unchanged, but alter the transformed
contribution.  We fix them first; Section~\ref{sec:interpolation} will correct
the positive-node coefficients to impose all contacts of the coupled pair.
We choose the following values, whose suitability will be established by the
norm and sign estimates below.  In the order of list 1 followed by list 2,
they are
\begin{equation}\label{eq:extra-data}
 E(k)=
 \begin{pmatrix}.567\\.250\\-.0074\\0\\-.209\\.0096\end{pmatrix}
 +\frac1k
 \begin{pmatrix}-.976\\.408\\.0115\\1\\.938\\-.017\end{pmatrix}.
\end{equation}
Let $E_1,E_2$ be $3\times10$ matrices.  Their first columns are the first
and last three entries, respectively, of the first vector in
\eqref{eq:extra-data}; their sixth columns use the corresponding entries of
the second vector, and all other columns vanish.  Since $v_1=1$, the stack
of $E_1q(k)$ and $E_2q(k)$ is exactly $E(k)$.

Define a $168\times10$ matrix $Y$ by the following nonzero entries, for
$1\le j\le5$:
\[
 Y_{d_{n_j},j}=-\frac1{Q_{n_j}},\qquad
 Y_{c_{n_j},j+5}=\frac1{Q_{n_j}},\qquad
 Y_{d_{n_j},j+5}=-\frac{D_{n_j}}{Q_{n_j}}.
\]
Then $Yq(k)$ is the jet projection of $T_k$ at these five nodes and is zero
at the other nodes of $\mathcal F$.  The remaining target jets are small;
Section~\ref{sec:interpolation} will include them in the correction.

To see which equations the finite columns must approximate, write
$A_1,A_2$ for their unknown positive-node parts, each a $168\times10$ matrix.
The coupled jet identities from Section~\ref{sec:fourier}, with integrals
replaced by quadrature and target data replaced by $Y$, give
\[
 A_1-DA_2=Y+\widetilde S_{\mathcal F e}E_2,
 \qquad
 A_2-DA_1=\widetilde S_{\mathcal F e}E_1.
\]
Each right side contains the other list's extras because the transformed
function in $H_i$ comes from that other list.  Denote these right sides by
$g_1,g_2$.  Adding and subtracting the equations gives
\[
 (I-D)(A_1+A_2)=g_1+g_2,
 \qquad
 (I+D)(A_1-A_2)=g_1-g_2.
\]
We use finite geometric sums in place of the two inverses and define
\begin{equation}\label{eq:finite-columns}
\begin{aligned}
 g_1&=Y+\widetilde S_{\mathcal F e}E_2,&
 g_2&=\widetilde S_{\mathcal F e}E_1,\\
 V_+&=\sum_{j=0}^{511}D^j,&
 V_-&=\sum_{j=0}^{511}(-D)^j,\\
 X_{1\mathcal F}
 &=\tfrac12\bigl(V_+(g_1+g_2)+V_-(g_1-g_2)\bigr),\\
 X_{2\mathcal F}
 &=\tfrac12\bigl(V_+(g_1+g_2)-V_-(g_1-g_2)\bigr),&
 X_i&=\begin{pmatrix}X_{i\mathcal F}\\E_i\end{pmatrix}\quad(i=1,2).
\end{aligned}
\end{equation}
The identities $(I-D)V_+=(I+D)V_-=I-D^{512}$ show that the error in
these equations is controlled by $D^{512}$.  Its bound, and the passage
back to the exact integrals, are proved in Section~\ref{sec:interpolation}.
Each $X_i$ has $171$ rows and ten columns; the reference list for a given
$k$ will be $X_iq(k)$.  A column is extended by zero beyond $M$ when
regarded as a full list.  For any such matrix $X$, the row
functions $p_X$ and $K_X$ have $j$th components $p_{(X)_j}$ and
$K_{p_{(X)_j}}$; the same convention applies to finite sums and derivatives.
Although the reference lists have finite support, their profiles $p_{X_iq(k)}$
and $K_{p_{X_iq(k)}}$ use the exact spectral integrals.  The tilde profiles
replace those integrals by the quadrature rule when constructing finite
moment rows.

The remaining finite bounds have three roles.  Bounds on powers of $D$ and
the sums $V_+,V_-$ will control the finite inverse in
Section~\ref{sec:interpolation}, while column norms will control the size of
the reference lists and their interpolation residual.  Half-gap and tail
inequalities formed from the reference columns will supply quantitative
margins in the proof of $T_k-H_1\ge0$ and $H_2\ge0$.
Section~\ref{sec:signs} will compare these inequalities with the exact pair,
accounting for the interpolation correction and the approximation errors.
We next construct the rows for those sign inequalities.

\subsection{Containing parameter boxes}

The data vector $q(k)$ lies on a one-dimensional curve.  To obtain lower
bounds uniform in $k$, we place each part of that curve in a rectangular box.
The endpoints are
\[
 2.36,\quad2.65,\quad3.2,\quad4,\quad6,\quad\infty.
\]
For a real row $l=(a_1,\ldots,a_5,b_1,\ldots,b_5)$ and consecutive
endpoints $K,J$, define
\begin{equation}\label{eq:box-lower}
 [l]_{K,J}=\min_{x\in\{1/K,1/J\}}\left(
 a_1+xb_1+\sum_{j=2}^5
 \min_{H\in\{K,J\}}\{(a_j+xb_j)e^{-H(n_j-1)}\}\right).
\end{equation}
At $J=\infty$, the conventions are $1/J=0$ and
$e^{-J(n_j-1)}=0$ for $j>1$.

\begin{lemma}[Containing parameter boxes]\label{lem:parameter-box}
Let $K,J$ be consecutive endpoints in the displayed list.  For every real
row $l$ as above and every finite $k$ with $K\le k\le J$,
where the upper inequality is omitted if $J=\infty$, one has
$lq(k)\ge[l]_{K,J}$.
\end{lemma}
\begin{proof}
Write $x=1/k$ and $w_j=e^{-k(n_j-1)}$ for $j=2,\ldots,5$.
The physical vector lies in the box
\[
 x\in[1/J,1/K],\qquad
 w_j\in[e^{-J(n_j-1)},e^{-K(n_j-1)}].
\]
On this larger box the expression
$a_1+xb_1+\sum_{j=2}^5(a_j+xb_j)w_j$ is affine in each variable while
the others are fixed.  Its minimum therefore occurs at a vertex.  First
minimizing each $w_j$ and then the single shared $x$ gives exactly
\eqref{eq:box-lower}.  This is a containing-box bound, not sampling of the
curve $q(k)$.
\end{proof}

\subsection{Polynomial rows on the half-gaps}

We now construct the rows to which the containing-box bound will be applied.
They are Bernstein coefficients of polynomials that will control the signs
between nodes.  The required inequalities are $T_k-H_1\ge0$ and
$H_2\ge0$, so set $\sigma_1=-1$ and $\sigma_2=1$.
At each center
$m\in\mathcal N_0\cap[0,88]$, take the half-gap toward its successor and,
when $m>0$, the half-gap toward its predecessor in $\mathcal N_0$.  If $m'$
is that neighbor, put $y=(m'-m)/2$.  The half-gap consists of
$s=m+yx$, $0\le x\le1$, with signed $y$.

At a positive center the desired value and derivative contacts would make
$T_k-H_1$ and $H_2$ vanish to second order.  Before making an approximation,
we therefore subtract these two Taylor terms from each component.  For an
entire function $u$, define
\[
 \mathcal D_{m,2}u(s)
 =\frac{u(s)-u(m)-(s-m)u'(m)}{(s-m)^2},
 \qquad \mathcal D_{m,2}u(m)=\frac{u''(m)}2.
\]
If $u(m)=u'(m)=0$, then
\[
 u(s)=(s-m)^2\mathcal D_{m,2}u(s).
\]
Once exact interpolation is imposed, this identity applies to both
$u=T_k-H_1$ and $u=H_2$ at every positive center.  It reduces their sign
inequalities to lower bounds for quotients that remain regular at the nodes.
At the origin we use $\mathcal D_{0,0}u=u$.  Put $\nu=0$ there and $\nu=2$
at every positive center.  Taylor expansion gives
\[
 \mathcal D_{m,\nu}u(m+yx)
 =\sum_{r=0}^\infty\frac{y^r}{(r+\nu)!}u^{(r+\nu)}(m)x^r.
\]
This subtraction is defined even for reference functions that do not
interpolate exactly.  For the reference pair formed from the columns $X_i$,
we keep the first 41 coefficients and replace each spectral moment by its
quadrature sum.  The resulting ten-component rows are
\begin{equation}\label{eq:bernstein-rows}
\begin{aligned}
 l_{i,r}&=\sigma_i\frac{y^r}{(r+\nu)!}
 \left(\widetilde p_{X_i}^{(r+\nu)}(m)
      +\widetilde K_{X_{3-i}}^{(r+\nu)}(m)\right),&&0\le r\le40,\\
 L_{i,d}&=\sum_{r=0}^d\omega_{dr}l_{i,r},&
 \omega_{dr}&=\frac{\binom dr}{\binom{40}r},\qquad0\le d\le40.
\end{aligned}
\end{equation}
Equivalently, the first row is
\[
 l_{i,r}=\sigma_i y^{-\nu}\Rea\bigl(
 F_{r+\nu}(m,y)WX_i+G_{r+\nu}(m,y)WX_{3-i}\bigr).
\]
There is no singular case here: every selected half-gap has $y\ne0$.

The rows $L_{i,d}$ are exactly the degree-$40$ Bernstein coefficients of
$\sum_{r=0}^{40}l_{i,r}x^r$.  Indeed, the binomial theorem gives
\[
 x^r=\sum_{d=r}^{40}\frac{\binom dr}{\binom{40}r}
       \binom{40}d x^d(1-x)^{40-d}.
\]
On $[0,1]$ the complete family
$\binom{40}d x^d(1-x)^{40-d}$, $0\le d\le40$, is nonnegative and sums
to one.  Thus a polynomial value lies between its smallest and largest
Bernstein coefficients, the standard coefficient-range enclosure discussed
by Titi and Garloff~\citep[Section~3]{TitiGarloff2019}.  The identity just proved specifies
our rows; the local remainder and parameter estimates are separate.

The first inequality also contains the target quotient
$\mathcal D_{m,\nu}T_k$.  To bound it below on the same half-gap and
parameter box, define
\begin{equation}\label{eq:target-support}
 (a_r^*)_{r=0}^{40}=
 \begin{cases}
 \bigl((e^K/K)(-Ky)^r/r!\bigr)_{r=0}^{40},&m=0,\\
 \bigl(K(-Ky)^r/(r+2)!\bigr)_{r=0}^{40},&m=1,\\
 \bigl((J/2)e^{-13J/6},0,\ldots,0\bigr),&m=3,\ J<\infty,\\
 (0,\ldots,0),&\text{otherwise}.
 \end{cases}
\end{equation}
Lemma~\ref{lem:exponential-support} proves the precise lower bound
\[
 \mathcal D_{m,\nu}T_k(m+yx)
 \ge\sum_{r=0}^{40}a_r^*x^r-.000001\qquad(0\le x\le1).
\]
Thus the target polynomial supports the same quotient as the reference
polynomial; its truncation error is retained separately.

\subsection{The complete finite certificate}

For large $s$, we will use the low-node terms of the cardinal formula
rather than its Taylor expansion.  Put
\[
 s_0=89.5,\qquad \mathcal N_{\le88}=\mathcal N_0\cap[0,88].
\]
Write $u_{i,n}$ and $d^*_{i,n}$ for the $c_n$ and $d_n$ rows of
$\sigma_iX_i$, and $C_i^*$ for its constant row.  For the reference list
$X_iq(k)$, the signed rational part from these rows is
\[
 \sigma_iR_{i,0}^{\mathrm{ref}}(s)
 =C_i^*q(k)+\sum_{n\in\mathcal N_{\le88}}
 \left(\frac{u_{i,n}q(k)}{(s-n)^2}
       +\frac{d^*_{i,n}q(k)}{s-n}\right).
\]
The identity
\[
 \frac1{s-n}=\frac1s+\frac{n}{s(s-n)}
\]
groups the simple-pole rows into a coefficient of $1/s$ and a remaining
sum.  For a fixed parameter box, put
$\eta_{K,J}(l)=\min(0,[l]_{K,J})$.  This gives a nonpositive lower bound
for the numerator $lq(k)$.  On $s\ge s_0$, the reciprocal factors are
positive and decrease with $s$, so their products with this lower bound
are smallest at $s_0$.  Keeping the full lower bound for the constant row
gives the tail expression in the proposition below.  In
Section~\ref{sec:signs}, its positive margin will absorb the correction
to the exact lists and the omitted high-node and transformed terms.

\begin{proposition}[Finite certificate]\label{prop:finite-certificate}
The exact arrays defined in this section satisfy
\begin{equation}\label{eq:finite-norms}
\begin{gathered}
 \norm{D^{64}}<.004,\qquad \norm{V_+}+\norm{V_-}<36,\\
 \sum_{i=1}^2\norm{(X_i)_j}_1<
 \begin{cases}2.3,&1\le j\le5,\\5.6,&6\le j\le10.\end{cases}
\end{gathered}
\end{equation}
For every specified half-gap, $i=1,2$, consecutive parameter endpoints
$K,J$, and $0\le d\le40$,
\begin{equation}\label{eq:bernstein-certificate}
 [L_{i,d}]_{K,J}
 +\mathbf1_{i=1}\sum_{r=0}^d\omega_{dr}a_r^*>.001.
\end{equation}
For every box and $i=1,2$,
\begin{equation}\label{eq:tail-certificate}
\begin{split}
 [C_i^*]_{K,J}
 &+\frac{\eta_{K,J}\bigl(\sum_{n\in\mathcal N_{\le88}}d^*_{i,n}\bigr)}{s_0}\\
 &+\sum_{n\in\mathcal N_{\le88}}\left(
 \frac{\eta_{K,J}(u_{i,n})}{(s_0-n)^2}
 +\frac{n\,\eta_{K,J}(d^*_{i,n})}{s_0(s_0-n)}\right)>.0021.
\end{split}
\end{equation}
\end{proposition}

\begin{proof}[Finite verification]
We specify the complete stronger exact-array comparisons in
Tables~\ref{table:power-bounds}--\ref{table:sign-bounds} and the route by
which each is checked.  Their looser consequences are the proposition's
outputs, whose analytic use is separate.  No precomputed coefficient list or approximate auxiliary function
is an input: all arrays are reconstructed from \eqref{eq:quadrature-rule}--
\eqref{eq:finite-columns}.

The power sums use nine doubling steps.  Start with $U=D$ and $V_+=V_-=I$.
At step $r=0,\ldots,8$, replace $V_+$ by $V_+(I+U)$ and $V_-$ by
$V_-(I+\epsilon_rU)$, where $\epsilon_0=-1$ and $\epsilon_r=1$ for $r>0$;
square $U$ before the next step.  Before step $r$, $U=D^{2^r}$, and
\[
 \prod_{r=0}^8(I+D^{2^r})=\sum_{j=0}^{511}D^j,
 \qquad
 (I-D)\prod_{r=1}^8(I+D^{2^r})=\sum_{j=0}^{511}(-D)^j.
\]
Thus the final sums are exactly those in \eqref{eq:finite-columns}.
Table~\ref{table:power-bounds} gives strict upper bounds for $U$ before each
update and for both sums after it.  Row 6 therefore bounds $D^{64}$.

\begin{table}[tb]
\centering
\caption{Strict upper bounds in the nine-step power computation.}
\label{table:power-bounds}
\begin{tabular}{@{}rrrr@{}}
\toprule
Step $r$&$\norm{U}$&$\norm{V_+}$&$\norm{V_-}$\\
\midrule
0&1.74&2.74&2.1\\
1&2.09&6.94&2.2\\
2&2.27&15.32&2.13\\
3&1.75&25.3&1.97\\
4&.75&31.1&1.88\\
5&.12&32.09&1.85\\
6&.0031&32.12&1.85\\
7&.000002&32.12&1.85\\
8&.000000000001&32.12&1.85\\
\bottomrule
\end{tabular}
\end{table}

\begin{table}[tb]
\centering
\caption{Strict upper bounds for column norms.  In the last row the two
norms are added for the same column before taking the maximum.}
\label{table:column-bounds}
\begin{tabular}{@{}lrr@{}}
\toprule
Quantity&Columns $1$--$5$&Columns $6$--$10$\\
\midrule
$\max_j\norm{(X_1)_j}_1$&1.29&2.65\\
$\max_j\norm{(X_2)_j}_1$&1.01&3.03\\
$\max_j\sum_i\norm{(X_i)_j}_1$&2.20&5.40\\
\bottomrule
\end{tabular}
\end{table}

\begin{table}[tb]
\centering
\caption{Strict lower bounds for the left sides of
\eqref{eq:bernstein-certificate} and \eqref{eq:tail-certificate}.
Each half-gap entry covers both functions, every indicated center and signed
half-gap, and all Bernstein indices.  The right endpoint $J$ follows $K$ in
the endpoint list.}
\label{table:sign-bounds}
\begin{tabular}{@{}rrrrrr@{}}
\toprule
&\multicolumn{4}{c}{Half-gap certificate}&\\
\cmidrule(lr){2-5}
$K$&$m=0$&$m=1$&$m=3,4$&$7\le m\le88$&Tail\\
\midrule
2.36&.33&.019&.0033&.00116&.0022\\
2.65&.34&.017&.0037&.0018&.0029\\
3.2&.35&.033&.0061&.0026&.0036\\
4&.36&.036&.0085&.0031&.0043\\
6&.38&.011&.0105&.0037&.0052\\
\bottomrule
\end{tabular}
\end{table}

There are $46$ centers: the seven complete residue blocks in $[0,83]$
contribute $42$, and $84,85,87,88$ contribute four more.  Every center has
a right half-gap and every center except zero has a left one, giving $91$.
Consequently the sign table consists of $91\cdot2\cdot5\cdot41=37310$
separate real lower comparisons, and the tail column consists of ten.
The power and column tests likewise compare every indicated norm, including
each same-column paired norm.

The selected finite route encloses the exact mathematical quadrature arrays
directly by Arb ball arithmetic, whose inclusion principle is described by
Johansson~\citep{Johansson2017}.  The local source
\path{verification/numeric_balls.py} implements these formulas at $256$-bit
precision.  For an upper comparison the required test is that the whole ball
lies below its table cutoff; for a lower comparison it is that the whole ball
lies above it.  It performs these tests for all indices just counted.
All comparisons in Tables~\ref{table:power-bounds}--\ref{table:sign-bounds}
pass these strict inclusion tests.
The containing-box lemma, Bernstein range property, and analytic remainders
provide continuum coverage in $s$ and $k$.

Finally, the stronger cutoffs imply the displayed proposition bounds:
$.0031<.004$, $32.12+1.85<36$, $2.20<2.3$, $5.40<5.6$,
$.00116>.001$, and $.0022>.0021$.  Appendix~\ref{app:floor-certificate} gives a
second, fully specified floor-rounded rational procedure and proves that a
successful comparison with these same stronger cutoffs implies the looser
proposition bounds after its roundoff allowances.
\end{proof}
